% Figure: fatigue-life prediction. An S-N curve (stress amplitude versus
% cycles to failure, log-log straight line) on the left; a Miner's-rule
% damage accumulation, summing the damage each stress block contributes,
% on the right.
\begin{figure}[htbp]
\centering
\begin{tikzpicture}[font=\small]
% ===== Left panel: S-N curve =====
\begin{scope}[xshift=0cm]
\draw[->, thick, draw=engblack!70] (0,0) -- (5.2, 0)
  node[right, font=\scriptsize] {$\log_{10} N$ (cycles)};
\draw[->, thick, draw=engblack!70] (0,0) -- (0, 4.3)
  node[above, font=\scriptsize] {$\log_{10} S$ (stress amp.)};
\node[font=\small, anchor=south west] at (0.2, 4.1) {\textbf{(a) S-N curve}};
% Basquin straight line on log-log: log S = log A - (1/m) log N.
\draw[very thick, draw=engblue!80] (0.3, 3.8) -- (4.8, 0.5);
% Mark two operating stress blocks and read their lives off the line.
\draw[dashed, draw=engred!70] (0, 3.0) -- (1.35, 3.0) -- (1.35, 0);
\node[font=\scriptsize, anchor=east, color=engred!85] at (-0.05, 3.0) {$S_1$};
\node[font=\scriptsize, anchor=north, color=engred!85] at (1.35,-0.09) {$N_1$};
\draw[dashed, draw=engamber!80] (0, 1.7) -- (3.05, 1.7) -- (3.05, 0);
\node[font=\scriptsize, anchor=east, color=engamber!90] at (-0.05, 1.7) {$S_2$};
\node[font=\scriptsize, anchor=north, color=engamber!90] at (3.05,-0.09) {$N_2$};
\node[anchor=west, font=\scriptsize, color=engblack!70] at (2.3, 3.3)
  {slope $-1/m$};
\end{scope}
% ===== Right panel: Miner damage accumulation =====
\begin{scope}[xshift=7.2cm]
\draw[->, thick, draw=engblack!70] (0,0) -- (5.2, 0)
  node[right, font=\scriptsize] {service (blocks)};
\draw[->, thick, draw=engblack!70] (0,0) -- (0, 4.3)
  node[above, font=\scriptsize] {cumulative damage $D$};
\node[font=\small, anchor=south west] at (0.2, 4.1) {\textbf{(b) Miner sum}};
% Failure line at D = 1.
\draw[dashed, draw=engblack!55] (0, 3.4) -- (5.0, 3.4);
\node[font=\scriptsize, anchor=south east, color=engblack!70] at (5.0, 3.4) {$D=1$: failure};
% Staircase of damage increments, each block adding n_i/N_i.
\draw[very thick, draw=engred!85]
  (0,0) -- (0.7,0.5) -- (1.4,0.85) -- (2.1,1.5) -- (2.8,1.9)
  -- (3.5,2.6) -- (4.2,3.0) -- (4.7,3.4);
\filldraw[engred!85] (4.7,3.4) circle (0.06);
\node[font=\scriptsize, anchor=west, color=engred!85] at (3.0, 1.6)
  {$D=\sum_i n_i/N_i$};
\end{scope}
\end{tikzpicture}
\caption[Fatigue-life prediction by the S-N curve and Miner's rule]{Predicting
when a cyclically loaded part fails. (a) The S-N curve is a power law, a
straight line on log-log axes (the Basquin form $S^m N = \text{const}$), whose
slope $-1/m$ is read from constant-amplitude fatigue tests: a lower stress
amplitude buys a longer life. Two operating stress blocks $S_1$ and $S_2$ read
off their allowable lives $N_1$ and $N_2$. (b) Miner's linear rule sums the
fractional damage $n_i/N_i$ each block of $n_i$ cycles contributes; the part is
predicted to fail when the accumulated damage $D$ reaches unity. The rule is
the model, the S-N line is its identified parameter set, and the scatter in
real fatigue life around $D=1$ is the structural error the report must carry.}
\label{fig:vol02ch18:fatigue-life}
\end{figure}
